\chapter*{Введение}\label{ch:intro}
\addcontentsline{toc}{chapter}{Введение}


В современную эпоху цифровизации мобильные приложения становятся неотъемлемой частью повседневной жизни. Смартфоны используют около 70\% населения планеты, а на их долю приходится более 85\% всех используемых мобильных устройств - что создает благоприятную среду для мобильных сервисов и супераппов в целом \cite{ref1}.


Одновременно мусульманское население мира стремительно растет и составляет значительную часть глобальной аудитории: по оценкам, к 2030 году число мусульман может превысить 2,2 млрд человек - примерно 30\% населения Земли \cite{ref2}.


В рамках повседневного потребления для данной аудитории принципиальное значение имеет соблюдение требований халяль. Под халяльными продуктами понимаются продукты питания и товары, соответствующие нормам исламского права (шариата), включая запрет на использование определенных ингредиентов (например, свинины, алкоголя и их производных), а также требования к процессам производства, хранения и логистики.При этом определение халяльности продукта в современных условиях не всегда является однозначной задачей. Состав многих товаров включает сложные пищевые добавки, ферменты и ароматизаторы, происхождение которых затруднительно установить без специализированной информации. Дополнительную сложность создает отсутствие единых международных стандартов халяль и различия в требованиях сертифицирующих организаций, что осложняет осознанный выбор для конечного потребителя.Демографический сдвиг приводит к увеличению потребительских расходов в халяль-сегменте экономики. Согласно одной из оценок, общий объем глобального рынка халяльных продуктов питания и услуг к середине 2020-х годов исчисляется триллионами долларов - например, объем рынка халяльной еды превысил 2,6 трлн USD в 2024 году, и ожидается, что он будет расти более чем на 10\% в год \cite{ref3}.


На уровне повседневного потребления цифровые технологии выступают для мусульманских пользователей основным каналом доступа к информации и сервисам. Отраслевые аналитические обзоры указывают на высокий уровень мобильной активности данной аудитории: в ряде ключевых регионов проникновение смартфонов превышает 80\%, а мобильные устройства активно используются для получения контента, совершения покупок и взаимодействия с цифровыми услугами \cite{ref6}.


На цифровом рынке мобильных приложений в исламском пространстве уже есть яркие примеры больших и активно используемых продуктов. Так, популярное приложение Muslim Pro\textit{ }было скачано более 30 млн раз, а во время рамадана им пользовались от 3 до 4 млн активных пользователей в день \cite{ref4}. Кроме него существуют специализированные сервисы, такие как Muzz - мусульманское приложение для знакомств, которое насчитывает свыше 8 млн пользователей и сыграло роль в более чем 600 000 браков \cite{ref5}.


Вместе с тем существуют серьезные риски и ограничения, которые отражают религиозные и юридические рамки:

\begin{enumerate}
	\item отсутствие единых международных стандартов халяль, поскольку каждая страна и сертифицирующая организация формируют собственные правила - это усложняет проверку и соответствие продуктов и сервисов на глобальном рынке \cite{ref7}.
	\item правовые различия в отдельных юрисдикциях могут ограничивать распространение и функциональность приложений, если они затрагивают религиозные аспекты, например, интерпретацию исламских норм (фетвы), что приводит к необходимости локализации и соблюдения нормативов в каждой стране. (общая проблема стандартизации характерна для них как цифровых, так и физических халяль-продуктов) \cite{ref7}.
	\item религиозно-культурные требования требуют от цифровых продуктов высокой точности, толерантности к различным толкованиям и возможности кастомизации под локальные конфессии и подходы, что создает вызовы для UX/UI и алгоритмов рекомендаций.
\end{enumerate}

Эти факторы усиливают необходимость интеграции ИИ-механизмов, которые могут автоматически учитывать сложные правила, языковые особенности и культурный контекст в анализе и выдаче рекомендаций. Экосистема совмещающая мобильное приложение, серверную инфраструктуру и модели искусственного интеллекта может обеспечить комплексное, адаптивное и соответствующее религиозным ожиданиям решение для глобального рынка.


Целью работы является разработка, реализация и тестирование многофункционального iOS приложения для мусульман со сканером продуктов, чатом с искусственным интеллектом на основе Корана, напоминаниями о намазах и с текстами сур и аятов для повседневных религиозных практик пользователей.


Для достижения указанной цели были поставлены следующие задачи:

\begin{enumerate}
	\item Провести анализ предметной области, рассмотреть существующие решения для мусульман и проанализировать конкурентов.
	\item Сформулировать функциональные и нефункциональные требования к разрабатываемому приложению на основе анализа существующих решений.
	\item Разработать общую архитектуру, включая схему взаимодействия iOS-клиента, серверной части, базы данных и RAG-модуля.
	\item Спроектировать структуру базы данных для хранения и обработки информации о халяльности продуктов, пользовательских настройках, а также текстов Корана.
	\item Разработать API для обеспечения взаимодействия между клиентской и серверной частями приложения..
	\item Реализовать LLM сервис с RAG-пайплайном для формирования ответов на пользовательские вопросы в чате на основе текстов Корана.
	\item Разработать iOS-приложение, учитывающее функциональные и нефункциональные требования к продукту.
	\item Осуществить интеграцию клиентской и серверной частей, а также провести тестирование работоспособности всех модулей приложения и оценку качества ответов от RAG-системы.
\end{enumerate}

